\documentclass[12pt, letterpaper]{article}
\usepackage[utf8]{inputenc}
\usepackage[spanish]{babel}
\usepackage{amsmath, amssymb}
\usepackage{geometry}
\usepackage{graphicx}
\usepackage{xcolor}
\usepackage{enumitem}
\usepackage{array}
\usepackage{booktabs}
\usepackage{tcolorbox}
\tcbuselibrary{theorems, skins, breakable}

\geometry{margin=2.5cm}

\title{Fracciones Parciales}
\author{Jose Luciano Mejia Arias}
\date{}

\makeatletter
\renewcommand{\maketitle}{
  \begin{flushright}
    \includegraphics[width=3cm]{../logo_unal.png}
  \end{flushright}
  \begin{center}
    {\LARGE \textbf{\@title}}\\[0.3cm]
    {\large \@author}
  \end{center}
}
\makeatother

% ── Paleta de colores ──────────────────────────────────────────────────────────
\definecolor{azul}{RGB}{24, 95, 165}
\definecolor{azulclaro}{RGB}{230, 241, 251}
\definecolor{azultitulo}{RGB}{12, 68, 124}

\definecolor{verde}{RGB}{15, 110, 86}
\definecolor{verdeclaro}{RGB}{225, 245, 238}
\definecolor{verdetitulo}{RGB}{8, 80, 65}

\definecolor{ambar}{RGB}{133, 79, 11}
\definecolor{ambarclaro}{RGB}{250, 238, 218}
\definecolor{ambartitulo}{RGB}{99, 56, 6}

\definecolor{rosa}{RGB}{153, 53, 86}
\definecolor{rosaclaro}{RGB}{251, 234, 240}
\definecolor{rosatitulo}{RGB}{114, 36, 62}

% Colores nuevos (reemplazan los grises)
\definecolor{purpura}{RGB}{83, 74, 183}
\definecolor{purpuraclaro}{RGB}{238, 237, 254}
\definecolor{purpuratitulo}{RGB}{60, 52, 137}

\definecolor{coral}{RGB}{152, 60, 29}
\definecolor{coralclaro}{RGB}{250, 236, 231}
\definecolor{coraltitulo}{RGB}{113, 43, 19}

% ── Estilo base de cajas ───────────────────────────────────────────────────────
\tcbset{
  casobase/.style={
    enhanced, breakable,
    colback=white,
    boxrule=0.5pt,
    arc=7pt,
    left=12pt, right=12pt, top=10pt, bottom=10pt,
    fonttitle=\bfseries\normalsize,
    attach boxed title to top left={xshift=14pt, yshift=-\tcboxedtitleheight/2},
    boxed title style={arc=4pt, boxrule=0.5pt, left=8pt, right=8pt}
  },
  subcuadro/.style={
    enhanced, breakable,
    boxrule=0.4pt,
    arc=5pt,
    left=9pt, right=9pt, top=7pt, bottom=7pt,
    fonttitle=\small\bfseries,
    attach boxed title to top left={xshift=10pt, yshift=-\tcboxedtitleheight/2},
    boxed title style={arc=3pt, boxrule=0.4pt, left=6pt, right=6pt}
  }
}

% Cajas de casos
\newenvironment{caso1}{%
  \begin{tcolorbox}[casobase, colframe=azul, colbacktitle=azulclaro, coltitle=azultitulo,
    title=Caso 1 \textnormal{---} Factores lineales distintos]
}{\end{tcolorbox}}

\newenvironment{caso2}{%
  \begin{tcolorbox}[casobase, colframe=verde, colbacktitle=verdeclaro, coltitle=verdetitulo,
    title=Caso 2 \textnormal{---} Factores lineales repetidos]
}{\end{tcolorbox}}

\newenvironment{caso3}{%
  \begin{tcolorbox}[casobase, colframe=ambar, colbacktitle=ambarclaro, coltitle=ambartitulo,
    title=Caso 3 \textnormal{---} Factores cuadr\'{a}ticos irreducibles distintos]
}{\end{tcolorbox}}

\newenvironment{caso4}{%
  \begin{tcolorbox}[casobase, colframe=rosa, colbacktitle=rosaclaro, coltitle=rosatitulo,
    title=Caso 4 \textnormal{---} Factores cuadr\'{a}ticos irreducibles repetidos]
}{\end{tcolorbox}}

% Caja "Condicion previa" -- purpura
\newenvironment{condicion}{%
  \begin{tcolorbox}[subcuadro, colback=purpuraclaro, colframe=purpura,
    coltitle=purpuratitulo, colbacktitle=purpuraclaro,
    title=Condici\'{o}n previa]
}{\end{tcolorbox}}

% Caja "Como hallar las constantes" -- coral
\newenvironment{metodocuadro}[1]{%
  \begin{tcolorbox}[subcuadro, colback=coralclaro, colframe=coral,
    coltitle=coraltitulo, colbacktitle=coralclaro,
    title={#1}]
}{\end{tcolorbox}}

% Caja "Procedimiento general" -- purpura
\newenvironment{procedimiento}{%
  \begin{tcolorbox}[subcuadro, colback=purpuraclaro, colframe=purpura,
    coltitle=purpuratitulo, colbacktitle=purpuraclaro,
    title=Procedimiento general para hallar las constantes]
}{\end{tcolorbox}}

% ── Documento ─────────────────────────────────────────────────────────────────
\begin{document}
\maketitle

\begin{center}
  {\large\color{purpura} Descomposici\'{o}n de fracciones racionales}\\[4pt]
  \textcolor{purpura}{\rule{0.5\textwidth}{0.6pt}}
\end{center}

\vspace{0.4cm}

\begin{condicion}
Sea $\dfrac{P(x)}{Q(x)}$ una fracci\'{o}n racional \textbf{propia}: $\deg P(x) < \deg Q(x)$.\\[4pt]
Si no fuera el caso, primero se realiza la \textbf{divisi\'{o}n de polinomios} y se trabaja
\'{u}nicamente con el residuo. La descomposici\'{o}n depende de los \textbf{factores del denominador} $Q(x)$.
\end{condicion}

\vspace{0.5cm}

%% ── CASO 1 ───────────────────────────────────────────────────────────────────
\begin{caso1}

\textbf{Forma del denominador:} factores lineales $(x - a_i)$ todos distintos entre s\'{i}.

\[
  \frac{P(x)}{(x-a_1)(x-a_2)\cdots(x-a_n)}
  \;=\; \frac{A_1}{x-a_1} + \frac{A_2}{x-a_2} + \cdots + \frac{A_n}{x-a_n}
\]

\begin{metodocuadro}{C\'{o}mo hallar las constantes --- Caso 1}
Se multiplican ambos lados por el denominador completo:
\[
  P(x) = A_1(x-a_2)\cdots(x-a_n) + A_2(x-a_1)(x-a_3)\cdots(x-a_n) + \cdots
\]
\textbf{M\'{e}todo de ra\'{i}ces (el m\'{a}s r\'{a}pido):} se sustituye cada ra\'{i}z $x = a_i$.
Al hacerlo, todos los factores $(x-a_j)$ con $j \neq i$ se anulan, dejando una ecuaci\'{o}n
con una sola inc\'{o}gnita $A_i$.
\end{metodocuadro}

\medskip
\textbf{Ejemplo resuelto:}
\[
  \frac{3x + 1}{(x-1)(x+2)} = \frac{A}{x-1} + \frac{B}{x+2}
\]
Multiplicando por $(x-1)(x+2)$:
\[
  3x+1 = A(x+2) + B(x-1)
\]
\begin{itemize}[leftmargin=2em]
  \item $x = 1$:\quad $4 = 3A \;\Rightarrow\; A = \dfrac{4}{3}$
  \item $x = -2$:\quad $-5 = -3B \;\Rightarrow\; B = \dfrac{5}{3}$
\end{itemize}
\[
  \therefore\quad \frac{3x+1}{(x-1)(x+2)} = \frac{4/3}{x-1} + \frac{5/3}{x+2}
\]

\end{caso1}

\vspace{0.4cm}

%% ── CASO 2 ───────────────────────────────────────────────────────────────────
\begin{caso2}

\textbf{Forma del denominador:} un factor lineal $(x-a)$ elevado a la potencia $n$.

\[
  \frac{P(x)}{(x-a)^n}
  \;=\; \frac{A_1}{x-a} + \frac{A_2}{(x-a)^2} + \cdots + \frac{A_n}{(x-a)^n}
\]

Si el denominador mezcla un factor repetido con uno simple:
\[
  \frac{P(x)}{(x-a)^n(x-b)}
  = \frac{A_1}{x-a} + \frac{A_2}{(x-a)^2} + \cdots + \frac{A_n}{(x-a)^n} + \frac{B}{x-b}
\]

\begin{metodocuadro}{C\'{o}mo hallar las constantes --- Caso 2}
Se multiplican ambos lados por todo el denominador.

\begin{itemize}[leftmargin=2em]
  \item \textbf{Constante $A_n$ (mayor potencia):} sustituir $x = a$. Todos los t\'{e}rminos
        salvo $A_n$ desaparecen, aislando directamente esa constante.
  \item \textbf{Constantes $A_1, \ldots, A_{n-1}$:} se obtienen \textbf{igualando coeficientes}
        de cada potencia de $x$ en ambos lados (sistema de ecuaciones), o bien
        diferenciando sucesivamente la identidad respecto a $x$ y evaluando en $x = a$.
  \item \textbf{Constante $B$} (factor simple $x-b$, si existe): sustituir $x = b$.
\end{itemize}
\end{metodocuadro}

\medskip
\textbf{Ejemplo resuelto:}
\[
  \frac{x^2 + 1}{(x-1)^2(x+3)} = \frac{A}{x-1} + \frac{B}{(x-1)^2} + \frac{C}{x+3}
\]
Multiplicando por $(x-1)^2(x+3)$:
\[
  x^2+1 = A(x-1)(x+3) + B(x+3) + C(x-1)^2
\]
\begin{itemize}[leftmargin=2em]
  \item $x = 1$:\quad $2 = 4B \;\Rightarrow\; B = \dfrac{1}{2}$
  \item $x = -3$:\quad $10 = 16C \;\Rightarrow\; C = \dfrac{5}{8}$
  \item Coeficiente de $x^2$:\quad $1 = A + C \;\Rightarrow\; A = \dfrac{3}{8}$
\end{itemize}

\end{caso2}

\vspace{0.4cm}

%% ── CASO 3 ───────────────────────────────────────────────────────────────────
\begin{caso3}

\textbf{Forma del denominador:} factores cuadr\'{a}ticos $ax^2+bx+c$ con $b^2-4ac < 0$
(sin ra\'{i}ces reales), todos distintos.

\[
  \frac{P(x)}{(a_1x^2+b_1x+c_1)\cdots(a_kx^2+b_kx+c_k)}
  \;=\; \sum_{i=1}^{k} \frac{A_i x + B_i}{a_i x^2 + b_i x + c_i}
\]

\begin{metodocuadro}{C\'{o}mo hallar las constantes --- Caso 3}
Los factores cuadr\'{a}ticos irreducibles \textbf{no tienen ra\'{i}ces reales}, por lo que
el m\'{e}todo de sustituci\'{o}n directa no puede aislar cada par $(A_i, B_i)$ por s\'{i} solo.
El procedimiento es:

\begin{enumerate}[leftmargin=2em, label=\arabic*.]
  \item Multiplicar ambos lados por el denominador completo.
  \item \textbf{Igualar coeficientes} de cada potencia de $x$ (grado 0, 1, 2, \ldots) en
        ambos lados. Esto genera un sistema de ecuaciones lineales con las inc\'{o}gnitas
        $A_1, B_1, A_2, B_2, \ldots$
  \item Resolver el sistema resultante.
\end{enumerate}

\textit{Nota:} si el denominador mezcla factores lineales y cuadr\'{a}ticos, las ra\'{i}ces
reales de los factores lineales pueden sustituirse primero para reducir el sistema.
\end{metodocuadro}

\medskip
\textbf{Ejemplo resuelto:}
\[
  \frac{5x}{(x^2+1)(x^2+4)} = \frac{Ax+B}{x^2+1} + \frac{Cx+D}{x^2+4}
\]
Multiplicando por $(x^2+1)(x^2+4)$:
\[
  5x = (Ax+B)(x^2+4) + (Cx+D)(x^2+1)
\]
Igualando coeficientes:
\[
  \begin{cases}
    x^3\!\!: & A + C = 0 \\
    x^2\!\!: & B + D = 0 \\
    x^1\!\!: & 4A + C = 5 \\
    x^0\!\!: & 4B + D = 0
  \end{cases}
  \;\Rightarrow\;
  A = \tfrac{5}{3},\; C = -\tfrac{5}{3},\; B = D = 0
\]
\[
  \therefore\quad \frac{5x}{(x^2+1)(x^2+4)} = \frac{5x/3}{x^2+1} - \frac{5x/3}{x^2+4}
\]

\end{caso3}

\vspace{0.4cm}

%% ── CASO 4 ───────────────────────────────────────────────────────────────────
\begin{caso4}

\textbf{Forma del denominador:} un factor cuadr\'{a}tico irreducible $(ax^2+bx+c)$ repetido $n$ veces.

\[
  \frac{P(x)}{(ax^2+bx+c)^n}
  \;=\; \frac{A_1 x+B_1}{ax^2+bx+c} + \frac{A_2 x+B_2}{(ax^2+bx+c)^2}
        + \cdots + \frac{A_n x+B_n}{(ax^2+bx+c)^n}
\]

\begin{metodocuadro}{C\'{o}mo hallar las constantes --- Caso 4}
Al igual que en el Caso 3, no existen ra\'{i}ces reales para sustituir directamente.
El \'{u}nico m\'{e}todo disponible es la \textbf{igualaci\'{o}n de coeficientes}:

\begin{enumerate}[leftmargin=2em, label=\arabic*.]
  \item Multiplicar ambos lados por $(ax^2+bx+c)^n$.
  \item Expandir el lado derecho completamente.
  \item Igualar los coeficientes de $x^0, x^1, x^2, \ldots$ de ambos lados.
        Esto produce un sistema de $2n$ ecuaciones con $2n$ inc\'{o}gnitas
        $(A_1, B_1, \ldots, A_n, B_n)$.
  \item Resolver el sistema por sustituci\'{o}n o eliminaci\'{o}n.
\end{enumerate}

\textit{Nota:} si el denominador combina un factor cuadr\'{a}tico repetido con factores lineales,
primero se sustituyen las ra\'{i}ces reales para hallar esas constantes y luego se reduce el sistema.
\end{metodocuadro}

\medskip
\textbf{Ejemplo resuelto:}
\[
  \frac{x^3+2x}{(x^2+1)^2} = \frac{Ax+B}{x^2+1} + \frac{Cx+D}{(x^2+1)^2}
\]
Multiplicando por $(x^2+1)^2$:
\[
  x^3+2x = (Ax+B)(x^2+1) + (Cx+D)
\]
Expandiendo: $\;x^3+2x = Ax^3 + Bx^2 + (A+C)x + (B+D)$.\\[4pt]
Igualando coeficientes:
\[
  \begin{cases}
    x^3\!\!: & A = 1 \\
    x^2\!\!: & B = 0 \\
    x^1\!\!: & A + C = 2 \;\Rightarrow\; C = 1 \\
    x^0\!\!: & B + D = 0 \;\Rightarrow\; D = 0
  \end{cases}
\]
\[
  \therefore\quad \frac{x^3+2x}{(x^2+1)^2} = \frac{x}{x^2+1} + \frac{x}{(x^2+1)^2}
\]

\end{caso4}

\vspace{0.5cm}

%% ── PROCEDIMIENTO GENERAL ────────────────────────────────────────────────────
\begin{procedimiento}
\begin{enumerate}[leftmargin=*, label=\textbf{\arabic*.}]
  \item \textbf{Verificar el grado.} $\deg P < \deg Q$. Si no, dividir primero y trabajar con el residuo.
  \item \textbf{Factorizar} $Q(x)$ en factores lineales y cuadr\'{a}ticos irreducibles.
  \item \textbf{Escribir la forma} de las fracciones parciales seg\'{u}n los casos 1--4.
  \item \textbf{Multiplicar} ambos lados por $Q(x)$ para eliminar denominadores.
  \item \textbf{Hallar las constantes:}
    \begin{itemize}
      \item \textit{Ra\'{i}ces reales} (Casos 1 y 2): sustituir $x = a_i$ para aislar cada constante.
      \item \textit{Factores cuadr\'{a}ticos} (Casos 3 y 4): igualar coeficientes y resolver el sistema.
      \item \textit{Casos mixtos}: combinar ambos m\'{e}todos.
    \end{itemize}
  \item \textbf{Escribir el resultado final} con las constantes ya determinadas.
\end{enumerate}
\end{procedimiento}

\vspace{0.5cm}

%% ── TABLA RESUMEN ────────────────────────────────────────────────────────────
\begin{center}
{\bfseries\large Tabla resumen}\\[10pt]
\renewcommand{\arraystretch}{2}
\begin{tabular}{>{\centering\arraybackslash}m{4.8cm}
                >{\centering\arraybackslash}m{5.5cm}
                >{\centering\arraybackslash}m{3.2cm}}
  \toprule
  \textbf{Factor en } $Q(x)$ & \textbf{Fracci\'{o}n(es) parcial(es)} & \textbf{M\'{e}todo} \\
  \midrule
  $(x-a)$ simple
    & $\dfrac{A}{x-a}$
    & Ra\'{i}ces \\
  $(x-a)^n$ repetido
    & $\dfrac{A_1}{x-a}+\cdots+\dfrac{A_n}{(x-a)^n}$
    & Ra\'{i}ces + coef. \\
  $ax^2{+}bx{+}c$ irreducible
    & $\dfrac{Ax+B}{ax^2+bx+c}$
    & Coeficientes \\
  $(ax^2{+}bx{+}c)^n$ repetido
    & $\dfrac{A_1x+B_1}{ax^2+bx+c}+\cdots+\dfrac{A_nx+B_n}{(ax^2+bx+c)^n}$
    & Coeficientes \\
  \bottomrule
\end{tabular}
\end{center}

\end{document}
